\documentclass[11pt,a4paper]{article}

\usepackage[margin=0.9in,headheight=14pt]{geometry}
\usepackage{amsmath,amssymb,amsthm,mathtools}
\usepackage{enumitem}
\usepackage{xcolor}
\usepackage{microtype}
\usepackage[hidelinks]{hyperref}
\usepackage{fancyhdr}

\definecolor{titleblue}{RGB}{24,65,115}
\definecolor{solutiongreen}{RGB}{20,90,60}

\setlength{\parindent}{0pt}
\setlength{\parskip}{0.48em}
\setlength{\emergencystretch}{2em}
\allowdisplaybreaks[2]

\pagestyle{fancy}
\fancyhf{}
\fancyhead[L]{MTL712}
\fancyhead[R]{Note 3 Practice}
\fancyfoot[C]{\thepage}

\newtheoremstyle{questionstyle}
  {10pt}{6pt}{\normalfont}{}{\bfseries\color{titleblue}}{.}{0.5em}{}
\theoremstyle{questionstyle}
\newtheorem{question}{Question}

\newenvironment{solution}
  {\par\noindent\textbf{\color{solutiongreen}Solution.}\ }
  {\hfill$\square$\par}

\newcommand{\pointmark}[1]{\hfill\textnormal{\textit{[#1 marks]}}}
\newcommand{\R}{\mathbb{R}}
\newcommand{\abs}[1]{\left|#1\right|}
\newcommand{\norm}[1]{\left\lVert#1\right\rVert}
\newcommand{\fall}[1]{n^{(#1)}}

\title{\textcolor{titleblue}{\textbf{MTL712: Computational Methods for Differential Equations}}\\[5pt]
\large Note 3 Practice Questions with Complete Solutions}
\author{}
\date{}

\begin{document}

\maketitle
\thispagestyle{plain}

\begin{center}
\emph{Root conditions, stability, order conditions, and difference equations}
\end{center}

\tableofcontents
\clearpage

% =====================================================================
\section{Proof-Based Questions}

\begin{question}[Necessity of the root condition]\pointmark{8}
Prove that a linear multistep method cannot be zero-stable if its first
characteristic polynomial \(\rho\) has either
\begin{enumerate}[label=(\alph*)]
  \item a root \(z\) with \(\abs{z}>1\), or
  \item a repeated root \(z\) with \(\abs{z}=1\).
\end{enumerate}
\end{question}

\begin{solution}
Apply the method to \(y'=0\).  The resulting homogeneous recurrence is
\[
  \sum_{j=0}^{k}\alpha_jy_{n+j}=0,
\]
and its auxiliary polynomial is \(\rho(z)\).  If \(z_s\) is a root of
multiplicity \(m_s\), the corresponding terms in the general solution
are
\[
  p_s(n)z_s^n,\qquad \deg p_s\leq m_s-1.
\]

If \(\abs{z_s}>1\), choose starting values that produce the solution
\(y_n=z_s^n\).  Its initial values are bounded independently of \(h\),
but, on a fixed interval \(0\leq nh\leq T\), the value at
\(n=\lfloor T/h\rfloor\) grows without bound as \(h\to0\).

If \(\abs{z_s}=1\) and \(m_s>1\), choose starting values producing
\(y_n=n^{m_s-1}z_s^n\).  This also becomes unbounded when
\(n=\lfloor T/h\rfloor\).

Compare either sequence with the zero solution.  No constant \(K\)
independent of \(h\) can satisfy
\[
  \abs{y_n-\widehat y_n}
  \leq K\max_{0\leq j<k}\abs{y_j-\widehat y_j}.
\]
Hence the method is not zero-stable.
\end{solution}

\begin{question}[Sufficiency of the root condition]\pointmark{10}
Assume that \(f\) is uniformly Lipschitz in \(y\) and that the root
condition holds.  Prove zero-stability using bounded powers of the
companion matrix, variation of constants, and a discrete Gronwall
inequality.
\end{question}

\begin{solution}
Write the homogeneous recurrence as a first-order system
\[
  Y_{n+1}=AY_n,
\]
where \(A\) is the companion matrix associated with \(\rho\).  The
eigenvalues of \(A\) are the roots of \(\rho\).  In Jordan form, blocks
whose eigenvalues lie strictly inside the unit disk decay, while a
block corresponding to a root on the unit circle is bounded precisely
when that root is simple.  The root condition therefore implies
\[
  \sup_{m\geq0}\norm{A^m}\leq M
\]
for some constant \(M\).

Let \(e_n=y_n-\widehat y_n\) be the difference of two numerical
solutions.  After writing the error recurrence in first-order form,
variation of constants gives
\[
  \mathcal E_n
  =A^{n-k+1}\mathcal E_{k-1}
   +h\sum_{\ell=k-1}^{n-1}
      A^{n-1-\ell}G_\ell ,
\]
where \(\mathcal E_n\) contains \(k\) consecutive errors.  If \(f\) is
Lipschitz in \(y\), then
\[
  \norm{G_\ell}\leq C\norm{\mathcal E_{\ell+1}}.
\]
For an implicit method, the term involving the newest error can be
absorbed into the left side for sufficiently small \(h\).  Consequently,
\[
  \norm{\mathcal E_n}
  \leq M\norm{\mathcal E_{k-1}}
    +Ch\sum_{\ell=k-1}^{n-1}\norm{\mathcal E_\ell}.
\]
The discrete Gronwall inequality yields
\[
  \norm{\mathcal E_n}
  \leq M e^{CT}\norm{\mathcal E_{k-1}},
  \qquad nh\leq T.
\]
The bound is independent of \(h\), which proves zero-stability.
\end{solution}

\begin{question}[Consistency from the local truncation error]\pointmark{8}
Assume \(\sigma(1)\neq0\).  Starting from the normalized local
truncation error, prove that a linear multistep method is consistent
if and only if
\[
  \rho(1)=0,\qquad \rho'(1)=\sigma(1).
\]
\end{question}

\begin{solution}
For
\[
  L[y;h]=\sum_{j=0}^{k}\alpha_jy(x+jh)
  -h\sum_{j=0}^{k}\beta_jy'(x+jh),
\]
Taylor expansion about \(x\) gives
\[
  L[y;h]=c_0y(x)+c_1hy'(x)+c_2h^2y''(x)+\cdots,
\]
where
\[
  c_0=\sum_{j=0}^{k}\alpha_j=\rho(1)
\]
and
\[
  c_1=\sum_{j=0}^{k}j\alpha_j-\sum_{j=0}^{k}\beta_j
      =\rho'(1)-\sigma(1).
\]
Assuming \(\sigma(1)\neq0\), the normalized local truncation error is
\[
  T_n=\frac1{\sigma(1)}
  \left(\frac{c_0}{h}y(x_n)+c_1y'(x_n)
  +c_2hy''(x_n)+\cdots\right).
\]
For \(T_n\to0\) along every sufficiently smooth solution, it is
necessary and sufficient that \(c_0=c_1=0\).  These equations are
exactly
\[
  \rho(1)=0,\qquad \rho'(1)=\sigma(1).
\]
\end{solution}

\begin{question}[Shift invariance of order]\pointmark{6}
Prove that replacing every index \(n+j\) by \(n+j+t\), where \(t\) is
a fixed integer, does not change the order or error constant of a
linear multistep method.
\end{question}

\begin{solution}
The shifted difference operator is
\[
  L_t[y(x);h]
  =\sum_{j=0}^{k}\alpha_jy(x+(j+t)h)
  -h\sum_{j=0}^{k}\beta_jy'(x+(j+t)h).
\]
It satisfies
\[
  L_t[y(x);h]=L[y(x+th);h].
\]
If \(L\) has order \(p\), then
\[
  L[y(x);h]
  =c_{p+1}h^{p+1}y^{(p+1)}(x)+O(h^{p+2}).
\]
Therefore
\[
  L_t[y(x);h]
  =c_{p+1}h^{p+1}y^{(p+1)}(x+th)+O(h^{p+2})
  =c_{p+1}h^{p+1}y^{(p+1)}(x)+O(h^{p+2}).
\]
The first nonzero coefficient remains \(c_{p+1}\), while
\(\sigma(1)\) is unchanged.  Thus both the order and
\(c_{p+1}/\sigma(1)\) are unchanged.
\end{solution}

\begin{question}[Order characterization using \(\Psi\)]\pointmark{10}
Let \(p\geq1\), \(\rho(1)=0\), and \(\sigma(1)\neq0\), and interpret
\(\rho(z)/\log z\) by its removable extension at \(z=1\).  Prove that
a linear multistep method has order \(p\) if and only if
\[
  \Psi(z)=\frac{\rho(z)}{\log z}-\sigma(z)
\]
has a zero of multiplicity \(p\) at \(z=1\).
\end{question}

\begin{solution}
Choose a branch of \(\log z\) near \(z=1\).  Put \(z=e^s\) and define
\[
  g(s)=\rho(e^s)-s\sigma(e^s).
\]
At \(q=0\), \(g(0)=c_0\).  For \(q\geq1\), direct differentiation at
\(s=0\) gives
\[
  g^{(q)}(0)
  =\sum_{j=0}^{k}\left(j^q\alpha_j
  -qj^{q-1}\beta_j\right)=q!c_q.
\]
Thus an order-\(p\) method, for which
\[
  c_0=\cdots=c_p=0,\qquad c_{p+1}\neq0,
\]
makes \(g\) vanish to order \(p+1\) at \(s=0\).

Moreover,
\[
  \Psi(e^s)
  =\frac{\rho(e^s)-s\sigma(e^s)}{s}
  =\frac{g(s)}s.
\]
Hence \(\Psi(e^s)\) vanishes to order \(p\) at \(s=0\).  Since
\(s\mapsto e^s\) is locally invertible at \(s=0\), this is equivalent
to \(\Psi(z)\) having a zero of multiplicity \(p\) at \(z=1\).
Reversing the argument proves the converse.
\end{solution}

\begin{question}[Construction of the second characteristic polynomial]\pointmark{8}
Let \(\rho\) have degree \(k\), with
\[
  \rho(1)=0,\qquad \rho'(1)\neq0.
\]
Prove that there is a unique polynomial \(\sigma\), of degree at most
\(k\), for which the method has order at least \(k+1\).
\end{question}

\begin{solution}
The function
\[
  R(z)=\frac{\rho(z)}{\log z}
\]
has a removable singularity at \(z=1\), and
\[
  R(1)=\rho'(1).
\]
Let \(\sigma\) be the degree-\(k\) Taylor polynomial of \(R\) at \(1\):
\[
  \sigma(z)=\sum_{m=0}^{k}
  \frac{R^{(m)}(1)}{m!}(z-1)^m.
\]
Then \(R-\sigma\) has a zero of multiplicity at least \(k+1\) at
\(z=1\).  By the order characterization, the resulting method has
order at least \(k+1\).

If \(\widetilde\sigma\) were another such polynomial, then
\(\sigma-\widetilde\sigma\), a polynomial of degree at most \(k\),
would have a zero of multiplicity at least \(k+1\).  It must therefore
be identically zero.  This proves uniqueness.
\end{solution}

\begin{question}[Shift rule]\pointmark{6}
For \(\beta\neq0\), prove
\[
  \phi(E)\bigl[\beta^n v(n)\bigr]
  =\beta^n\phi(\beta E)v(n),
\]
and explain its usefulness when \(\phi(\beta)=0\).
\end{question}

\begin{solution}
For every integer \(j\geq0\),
\[
  E^j[\beta^nv(n)]
  =\beta^{n+j}v(n+j)
  =\beta^n(\beta E)^jv(n).
\]
Multiplying by the coefficients of \(\phi\) and adding proves the
identity.  Thus, if \(u\) is any particular solution of
\[
  \phi(\beta E)u(n)=v(n),
\]
then \(\beta^nu(n)\) is a particular solution of
\[
  \phi(E)y_n=\beta^nv(n).
\]
The symbolic inverse notation summarizes this rule, but solutions are
always determined only up to a complementary solution.
When \(\phi(\beta)=0\), the usual rule
\(\phi(E)^{-1}\beta^n=\beta^n/\phi(\beta)\) is undefined.  The shift
rule exposes factors of \(E-1=\Delta\), which can instead be inverted
using factorial polynomials.
\end{solution}

% =====================================================================
\section{Stability Questions}

\begin{question}[Strong and weak stability]\pointmark{6}
For
\[
  \rho(z)=(z-1)(z^2-2az+1),
\]
classify the associated method as strongly stable, weakly stable, or
not zero-stable for every \(a\in\R\).
\end{question}

\begin{solution}
The two roots of the quadratic have product \(1\).

If \(\abs a<1\), write \(a=\cos\theta\), where
\(0<\theta<\pi\).  The roots are \(e^{\pm i\theta}\).  They are
distinct, lie on the unit circle, and differ from \(1\).  Together
with the simple root \(1\), they satisfy the root condition.  The
method is weakly stable because it has additional unit-modulus roots.

If \(a=1\), the quadratic is \((z-1)^2\), so \(z=1\) is repeated.  If
\(a=-1\), it is \((z+1)^2\), so \(z=-1\) is repeated.  Both cases fail
the root condition.

If \(\abs a>1\), the quadratic has two real reciprocal roots and one
has modulus greater than \(1\).  Hence the method is not zero-stable.
There is no value of \(a\) for which the method is strongly stable.
\end{solution}

\begin{question}[Two-step Adams--Bashforth stability]\pointmark{8}
For
\[
  y_{n+2}=y_{n+1}+\frac h2(3f_{n+1}-f_n),
\]
derive the stability polynomial and use the quadratic Schur criterion
to find its real interval of absolute stability.
\end{question}

\begin{solution}
Applying the method to \(y'=\lambda y\), with \(\mu=h\lambda\), gives
\[
  \pi(z,\mu)
  =z^2-\left(1+\frac{3\mu}{2}\right)z+\frac{\mu}{2}.
\]
For a real monic quadratic \(z^2+az+b\), the roots lie in the open
unit disk if and only if
\[
  \abs b<1,\qquad 1+a+b>0,\qquad 1-a+b>0.
\]
Here
\[
  a=-1-\frac{3\mu}{2},\qquad b=\frac{\mu}{2}.
\]
The three inequalities become
\[
  \abs\mu<2,\qquad -\mu>0,\qquad 2+2\mu>0.
\]
Their intersection is
\[
  \boxed{-1<\mu<0}.
\]
\end{solution}

\begin{question}[BDF2 using Routh--Hurwitz]\pointmark{10}
Apply
\[
  3y_{n+2}-4y_{n+1}+y_n=2hf_{n+2}
\]
to \(y'=\lambda y\).  Use the Cayley transformation and the
Routh--Hurwitz criterion to determine its interval of absolute
stability on the negative real axis.
\end{question}

\begin{solution}
The stability polynomial is
\[
  \pi(z,\mu)=(3-2\mu)z^2-4z+1.
\]
Use \(z=(1+w)/(1-w)\) and multiply by \((1-w)^2\).  The transformed
polynomial is
\[
  q(w)=(8-2\mu)w^2+(4-4\mu)w-2\mu.
\]
For a quadratic, Routh--Hurwitz requires all three coefficients to
have the same positive sign.  If \(\mu<0\), then
\[
  8-2\mu>0,\qquad 4-4\mu>0,\qquad -2\mu>0.
\]
Thus every negative real \(\mu\) is stable.  At \(\mu=0\), the
principal root is \(1\), so the endpoint is excluded.  Therefore
\[
  \boxed{-\infty<\mu<0}.
\]
\end{solution}

\begin{question}[Three-step Adams--Bashforth stability]\pointmark{10}
Determine the real interval of absolute stability of
\[
  y_{n+3}=y_{n+2}
  +\frac h{12}(23f_{n+2}-16f_{n+1}+5f_n).
\]
\end{question}

\begin{solution}
The stability polynomial is
\[
  \pi(z,\mu)
  =z^3-z^2-\frac{\mu}{12}(23z^2-16z+5).
\]
Under \(z=(1+w)/(1-w)\), multiplication by \((1-w)^3\) gives
\[
  q(w)
  =\left(2+\frac{11\mu}{3}\right)w^3
   +\left(4-\frac{2\mu}{3}\right)w^2
   +(2-2\mu)w-\mu.
\]
For \(-6/11<\mu<0\), all four coefficients are positive.  The
remaining cubic Routh--Hurwitz condition is
\[
\begin{split}
 &\left(4-\frac{2\mu}{3}\right)(2-2\mu)
 -\left(2+\frac{11\mu}{3}\right)(-\mu)\\
 &\hspace{2cm}
 =8-\frac{22\mu}{3}+5\mu^2>0,
\end{split}
\]
which certainly holds for \(\mu<0\).

At \(\mu=-6/11\), the leading coefficient vanishes; equivalently,
\(z=-1\) lies on the stability boundary.  For \(\mu<-6/11\), the
leading and constant coefficients have opposite signs, so the
Routh--Hurwitz conditions cannot hold.  Positive \(\mu\) also gives
opposite signs.  Hence
\[
  \boxed{-\frac6{11}<\mu<0}.
\]
\end{solution}

\begin{question}[Relative but not absolute stability]\pointmark{6}
Construct a consistent two-step method and a value of \(\mu\) for
which the principal root dominates the parasitic root, but at least
one root has modulus greater than \(1\).
\end{question}

\begin{solution}
Take
\[
  \rho(z)=z^2-z,\qquad
  \sigma(z)=\frac{1+7z}{8}.
\]
Since
\[
  \rho(1)=0,\qquad \rho'(1)=1=\sigma(1),
\]
the method
\[
  y_{n+2}-y_{n+1}
  =\frac h8(f_n+7f_{n+1})
\]
is consistent.  At \(\mu=1\),
\[
  \pi(z,1)=z^2-\frac{15}{8}z-\frac18,
\]
whose roots are
\[
  z_{1,2}=\frac{15\pm\sqrt{257}}{16}.
\]
Numerically, these are approximately \(1.9395\) and \(-0.0645\).
The root continuing from \(1\) is the first root, and it dominates the
parasitic root.  Thus the method is relatively stable at \(\mu=1\),
but it is not absolutely stable because \(\abs{z_1}>1\).
\end{solution}

\begin{question}[Cayley transformation]\pointmark{5}
Prove that
\[
  w=\frac{z-1}{z+1}
\]
maps \(\abs z<1\) onto \(\operatorname{Re}w<0\).  Find the images of
\(z=1,-1,i\).
\end{question}

\begin{solution}
Using the inverse map \(z=(1+w)/(1-w)\),
\[
  \abs z<1
  \iff \abs{1+w}<\abs{1-w}.
\]
Squaring both sides gives
\[
  1+2\operatorname{Re}w+\abs w^2
  <1-2\operatorname{Re}w+\abs w^2,
\]
which is equivalent to \(\operatorname{Re}w<0\).  Therefore the open
unit disk maps onto the open left half-plane.

The specified points map as follows:
\[
  1\mapsto0,\qquad -1\mapsto\infty,\qquad
  i\mapsto i.
\]
The point \(-1\) is the pole of the transformation.
\end{solution}

% =====================================================================
\section{Order and Error Constants}

\begin{question}[Derivation of Simpson's multistep formula]\pointmark{10}
Determine \(a,b,c\) so that
\[
  y_{n+2}-y_n=h(af_n+bf_{n+1}+cf_{n+2})
\]
has the maximum possible order.  Find its order, error constant, and
zero-stability classification.
\end{question}

\begin{solution}
Here
\[
  \alpha_0=-1,\quad\alpha_1=0,\quad\alpha_2=1,
  \qquad
  \beta_0=a,\quad\beta_1=b,\quad\beta_2=c.
\]
The conditions \(c_1=c_2=c_3=0\) are
\[
  a+b+c=2,\qquad b+2c=2,\qquad b+4c=\frac83.
\]
They give
\[
  a=\frac13,\qquad b=\frac43,\qquad c=\frac13.
\]
The condition \(c_4=0\) then holds automatically, while
\[
  c_5=-\frac1{90}\neq0.
\]
Thus the method is
\[
  y_{n+2}-y_n
  =\frac h3(f_n+4f_{n+1}+f_{n+2})
\]
and has order \(4\).  Since \(\sigma(1)=2\), its normalized error
constant is
\[
  \boxed{C=\frac{c_5}{\sigma(1)}=-\frac1{180}}.
\]
The first characteristic polynomial is
\[
  \rho(z)=z^2-1=(z-1)(z+1).
\]
Both roots are simple and lie on the unit circle.  Hence the method is
zero-stable but weakly stable.
\end{solution}

\begin{question}[Unexpected extra order]\pointmark{8}
Take \(\rho(z)=z^2-1\).  Construct \(\sigma\) by matching the
degree-two Taylor polynomial of \(\rho(z)/\log z\) about \(z=1\).
Show that the resulting method has order four.
\end{question}

\begin{solution}
Set \(x=z-1\).  Since
\[
  z^2-1=2x+x^2,\qquad
  \log z=x-\frac{x^2}{2}+\frac{x^3}{3}-\cdots,
\]
division gives
\[
  \frac{z^2-1}{\log z}
  =2+2x+\frac{x^2}{3}+O(x^3).
\]
Therefore
\[
\begin{split}
  \sigma(z)
  &=2+2(z-1)+\frac{(z-1)^2}{3}\\
  &=\frac13(z^2+4z+1).
\end{split}
\]
This is precisely Simpson's multistep formula from the preceding
question.  Direct order-condition calculation gives
\[
  c_0=c_1=c_2=c_3=c_4=0,\qquad c_5=-\frac1{90}.
\]
Hence its order is \(4\), one higher than the order \(3\) guaranteed
by the general degree-\(k\) construction theorem.
\end{solution}

\begin{question}[Order conditions for a three-step family]\pointmark{8}
For
\[
  y_{n+3}+ay_{n+2}+by_{n+1}+cy_n
  =h(df_{n+3}+ef_{n+2}+gf_{n+1}),
\]
write all algebraic conditions required for order three.  State the
conditions for consistency and zero-stability.
\end{question}

\begin{solution}
The coefficients, indexed from \(j=0\), are
\[
  (\alpha_0,\alpha_1,\alpha_2,\alpha_3)=(c,b,a,1),
\quad
  (\beta_0,\beta_1,\beta_2,\beta_3)=(0,g,e,d).
\]
Order at least three requires \(c_0=c_1=c_2=c_3=0\), namely
\[
\begin{aligned}
  1+a+b+c&=0,\\
  3+2a+b-(d+e+g)&=0,\\
  \frac{9+4a+b}{2}-(3d+2e+g)&=0,\\
  \frac{27+8a+b}{6}-\frac{9d+4e+g}{2}&=0.
\end{aligned}
\]
For the order to be exactly three, one additionally requires
\[
  c_4=
  \frac{81+16a+b}{24}
  -\frac{27d+8e+g}{6}\neq0.
\]

Consistency requires only the first two displayed equations.  For
zero-stability, every root of
\[
  \rho(z)=z^3+az^2+bz+c
\]
must lie in \(\abs z\leq1\), and every root on \(\abs z=1\) must be
simple.
\end{solution}

% =====================================================================
\section{Difference Equations}

\begin{question}[Repeated and complex roots]\pointmark{7}
Find the real complementary solution of
\[
  (E-2)^2(E^2-2E+2)y_n=0
\]
and describe the asymptotic growth of its components.
\end{question}

\begin{solution}
The root \(2\) has multiplicity two.  The quadratic roots are
\[
  1\pm i=\sqrt2\,e^{\pm i\pi/4}.
\]
Hence the real complementary solution is
\[
  \boxed{
  y_n=(A+Bn)2^n
  +(\sqrt2)^n
   \left(C\cos\frac{n\pi}{4}
         +D\sin\frac{n\pi}{4}\right)}.
\]
The terms \(A2^n\) and \(Bn2^n\) grow respectively like \(2^n\) and
\(n2^n\).  The oscillatory pair has envelope \((\sqrt2)^n\).
Unless their coefficients vanish, all these components are unbounded.
\end{solution}

\begin{question}[Resonant exponential forcing]\pointmark{10}
Solve
\[
  (E-2)^2(E+1)y_n=2^n(n+1)
\]
using the shift rule.
\end{question}

\begin{solution}
The complementary solution is
\[
  y_n^{(c)}=(C_1+C_2n)2^n+C_3(-1)^n.
\]
For a particular solution, put \(y_n^{(p)}=2^nv(n)\).  The shift rule
gives
\[
  (E-2)^2(E+1)[2^nv(n)]
  =2^n(2E-2)^2(2E+1)v(n).
\]
Since \(2E-2=2\Delta\), \(v\) must satisfy
\[
  4\Delta^2(2E+1)v(n)=n+1.
\]
Take
\[
  v(n)=\frac{n^3}{72}-\frac{n^2}{36}
      =\frac{n^2(n-2)}{72}.
\]
Direct substitution gives
\[
  4\Delta^2(2E+1)v(n)=n+1.
\]
Therefore
\[
  \boxed{
  y_n=(C_1+C_2n)2^n+C_3(-1)^n
      +\frac{2^nn^2(n-2)}{72}}.
\]
The direct exponential rule fails because \(\phi(2)=0\), and in fact
\(2\) is a double root.
\end{solution}

\begin{question}[Polynomial resonance]\pointmark{8}
Solve
\[
  \Delta^3y_n=n^2+2n+1
\]
using factorial polynomials.
\end{question}

\begin{solution}
Write \(n^{(r)}=n(n-1)\cdots(n-r+1)\).  In this
falling-factorial basis,
\[
  n^2+2n+1=n^{(2)}+3n^{(1)}+n^{(0)}.
\]
Using
\[
  \Delta^m n^{(p)}
  =p(p-1)\cdots(p-m+1)n^{(p-m)},
\]
a particular solution is
\[
  y_n^{(p)}
  =\frac{n^{(5)}}{60}
   +\frac{n^{(4)}}8
   +\frac{n^{(3)}}6.
\]
The complementary equation \(\Delta^3y_n=0\) has solution
\(C_1+C_2n+C_3n^2\).  Hence
\[
  \boxed{
  y_n=C_1+C_2n+C_3n^2
  +\frac{n^{(5)}}{60}
  +\frac{n^{(4)}}8
  +\frac{n^{(3)}}6}.
\]
\end{solution}

\begin{question}[Mixed forcing]\pointmark{10}
Find the complete solution of
\[
  (E-1)^2(E+2)y_n=3^n+n^2.
\]
\end{question}

\begin{solution}
The complementary solution is
\[
  y_n^{(c)}=C_1+C_2n+C_3(-2)^n.
\]
For \(3^n\),
\[
  \phi(3)=(3-1)^2(3+2)=20,
\]
so its particular solution is \(3^n/20\).

For the polynomial forcing, write
\[
  \phi(E)=\Delta^2(\Delta+3).
\]
Since
\[
  \frac1{\Delta+3}
  =\frac13\left(1-\frac{\Delta}{3}
  +\frac{\Delta^2}{9}-\cdots\right),
\]
and \(\Delta^3n^2=0\),
\[
  \frac1{\Delta+3}n^2
  =\frac13n^2-\frac19(2n+1)+\frac{2}{27}.
\]
In factorial form this is
\[
  \frac13n^{(2)}+\frac19n^{(1)}-\frac1{27}.
\]
Applying \(1/\Delta^2\) gives
\[
  y_{n,\mathrm{poly}}^{(p)}
  =\frac{n^{(4)}}{36}
   +\frac{n^{(3)}}{54}
   -\frac{n^{(2)}}{54}.
\]
Therefore
\[
  \boxed{
  y_n=C_1+C_2n+C_3(-2)^n+\frac{3^n}{20}
  +\frac{n^{(4)}}{36}
  +\frac{n^{(3)}}{54}
  -\frac{n^{(2)}}{54}}.
\]
\end{solution}

\begin{question}[Boundedness of homogeneous solutions]\pointmark{8}
Give necessary and sufficient conditions on the roots of
\(\phi(\xi)\) for every solution of
\[
  \phi(E)y_n=0
\]
to remain bounded as \(n\to\infty\).  Prove your answer.
\end{question}

\begin{solution}
The necessary and sufficient conditions are
\[
  \abs{\xi}\leq1\quad\text{for every root},
  \qquad
  \abs{\xi}=1\Longrightarrow \xi
  \text{ is simple}.
\]
Indeed, a nonzero root \(\xi\) of multiplicity \(m\) contributes
\[
  (C_0+C_1n+\cdots+C_{m-1}n^{m-1})\xi^n.
\]
If \(\abs{\xi}<1\), every such term is bounded and in fact tends to
zero despite its polynomial factor.  If \(\abs{\xi}=1\) and \(m=1\),
the term \(C_0\xi^n\) is bounded.  Therefore the stated conditions are
sufficient.

Conversely, if \(\abs{\xi}>1\), the solution \(\xi^n\) is unbounded.
If \(\abs{\xi}=1\) and \(m>1\), the solution
\(n^{m-1}\xi^n\) is unbounded.  A zero root contributes only a finite
initial transient (for example, \(E^my_n=0\) forces \(y_n=0\) for
\(n\geq m\)), so it cannot cause asymptotic unboundedness.  Thus both
conditions are necessary and sufficient.
\end{solution}

% =====================================================================
\section{Comprehensive Numerical Case Drill}

This section fills the remaining computational case-types arising from
Note 3.  Together with Questions 8--20, it covers boundary and repeated
roots, direct Schur and Routh tests, parameterized stability,
consistency and convergence, all homogeneous root patterns, and the
principal nonhomogeneous forcing and resonance cases.

\begin{question}[Parameterized root condition]\pointmark{6}
Classify the zero-stability of a three-step method whose first
characteristic polynomial is
\[
  \rho(z)=(z-1)(z-r)^2,\qquad r\in\R.
\]
State whether each zero-stable case is strong or weak.
\end{question}

\begin{solution}
The root \(1\) is always present.  If \(\abs r<1\), the repeated root
\(r\) lies strictly inside the unit disk, where repeated roots are
allowed.  The only unit-modulus root is the simple root \(1\), so the
method is strongly stable.

If \(r=1\), then \(1\) is a triple root.  If \(r=-1\), then \(-1\) is
a repeated unit root.  Both violate the root condition.  If
\(\abs r>1\), a root lies outside the unit disk.  Therefore
\[
  \boxed{\text{strongly stable for }\abs r<1;
  \quad\text{not zero-stable for }\abs r\geq1.}
\]
There is no weakly stable case.
\end{solution}

\begin{question}[Recursive Schur test for a cubic]\pointmark{8}
Without finding its roots, use Schur's recursive criterion to decide
whether
\[
  \phi(z)=z^3-\frac12z^2+\frac14z-\frac18
\]
has all its roots in the open unit disk.
\end{question}

\begin{solution}
The conjugate-reciprocal polynomial is
\[
  \widehat\phi(z)
  =-\frac18z^3+\frac14z^2-\frac12z+1.
\]
The first Schur inequality is \(1>1/8\).  The reduced polynomial is
\[
\begin{split}
  \phi_1(z)
  &=\frac{\phi(z)+\frac18\widehat\phi(z)}{z}\\
  &=\frac3{64}(21z^2-10z+4).
\end{split}
\]
Multiplication by a nonzero constant does not change the roots, so
consider \(\psi(z)=21z^2-10z+4\).  Its first Schur inequality is
\(21>4\), and its reduction is
\[
  \psi_1(z)
  =\frac{21\psi(z)-4\widehat\psi(z)}{z}
  =425z-170.
\]
For this linear polynomial, \(425>170\), so its single root lies in
the open unit disk.  Applying the criterion backwards shows that
\[
  \boxed{\phi\text{ is a Schur polynomial}.}
\]
\end{solution}

\begin{question}[Quartic Routh--Hurwitz parameter range]\pointmark{8}
Let
\[
\begin{split}
  \phi_a(z)={}&(z-1)^4+2(z-1)^3(z+1)
  +3(z-1)^2(z+1)^2\\
  &+4(z-1)(z+1)^3+a(z+1)^4.
\end{split}
\]
Use the Cayley transformation and the quartic Routh--Hurwitz
criterion to determine all \(a\in\R\) for which every zero of
\(\phi_a\) lies in \(\abs z<1\).
\end{question}

\begin{solution}
Under \(w=(z-1)/(z+1)\), division by \((z+1)^4\) gives
\[
  q_a(w)=w^4+2w^3+3w^2+4w+a.
\]
Its coefficients are
\[
  a_0=1,\quad a_1=2,\quad a_2=3,\quad a_3=4,\quad a_4=a.
\]
They must all be positive, so \(a>0\).  The nontrivial quartic
Hurwitz determinant is
\[
\begin{split}
  a_1a_2a_3-a_0a_3^2-a_1^2a_4
  &=2\cdot3\cdot4-1\cdot4^2-2^2a\\
  &=8-4a.
\end{split}
\]
It is positive precisely when \(a<2\).  Thus
\[
  \boxed{0<a<2}.
\]
\end{solution}

\begin{question}[Complete real stability interval of the
\(\theta\)-method]\pointmark{8}
For \(0\leq\theta\leq1\), consider
\[
  y_{n+1}=y_n+h\bigl[(1-\theta)f_n+\theta f_{n+1}\bigr].
\]
Find its interval of absolute stability on the negative real axis for
every \(\theta\).  Determine when it contains the entire negative real
axis.
\end{question}

\begin{solution}
For \(y'=\lambda y\), the amplification factor is
\[
  R(\mu)=\frac{1+(1-\theta)\mu}{1-\theta\mu},
  \qquad \mu=h\lambda.
\]
For real \(\mu<0\), the denominator is positive.  The condition
\(\abs{R(\mu)}<1\) is equivalent to
\[
\begin{split}
  &(1-\theta\mu)^2-
  \bigl(1+(1-\theta)\mu\bigr)^2>0\\
  &\qquad\Longleftrightarrow
  (-\mu)\bigl[2+(1-2\theta)\mu\bigr]>0.
\end{split}
\]
Hence
\[
  \boxed{
  \begin{cases}
  -\dfrac{2}{1-2\theta}<\mu<0,
       &0\leq\theta<\dfrac12,\\[6pt]
  -\infty<\mu<0,
       &\dfrac12\leq\theta\leq1.
  \end{cases}}
\]
Thus the whole negative real axis is included exactly when
\(\theta\geq1/2\).  The cases \(\theta=0,1/2,1\) give forward Euler,
trapezoidal, and backward Euler, respectively.
\end{solution}

\begin{question}[Consistency, zero-stability, and convergence]\pointmark{10}
For each method, determine its order, zero-stability type, and whether
it is convergent:
\[
\begin{array}{ll}
\text{(a)}&
 y_{n+2}-y_n=2hf_{n+1},\\[2pt]
\text{(b)}&
 y_{n+2}-3y_{n+1}+2y_n=hf_{n+1},\\[2pt]
\text{(c)}&
 y_{n+2}-\dfrac43y_{n+1}+\dfrac13y_n
 =\dfrac23hf_{n+2}.
\end{array}
\]
\end{question}

\begin{solution}
\textbf{(a)} Here
\[
  \rho(z)=z^2-1,\qquad \sigma(z)=2z.
\]
Since \(\rho(1)=0\) and \(\rho'(1)=2=\sigma(1)\), it is consistent.
The coefficients satisfy \(c_2=0\), \(c_3=1/3\), so the order is two.
The roots \(1,-1\) are simple; hence the method is weakly zero-stable
and therefore convergent.

\textbf{(b)} Here
\[
  \rho(z)=(z-1)(z-2),\qquad \sigma(z)=z.
\]
Although \(\rho(1)=0\), one has \(\rho'(1)=-1\neq1=\sigma(1)\).
Thus \(c_0=0\) but \(c_1=-2\neq0\): the method has order zero and is
inconsistent.  It also has a root \(2\) outside the unit disk and is
not zero-stable.  It is not convergent.

\textbf{(c)} Here
\[
  \rho(z)=z^2-\frac43z+\frac13
  =(z-1)\left(z-\frac13\right),
  \qquad \sigma(z)=\frac23z^2.
\]
It is consistent because
\(\rho'(1)=2/3=\sigma(1)\).  Also \(c_2=0\) and
\(c_3=-2/9\), so it has order two.  Its roots are \(1\) and \(1/3\);
therefore it is strongly zero-stable and convergent.
\end{solution}

\begin{question}[Leapfrog truncation error]\pointmark{7}
For
\[
  y_{n+2}-y_n=2hf_{n+1},
\]
compute the normalized local truncation error on the exact solution
\(y(x)=e^x\).  Confirm the order and error constant directly.
\end{question}

\begin{solution}
Here \(\sigma(1)=2\).  Substituting \(y=e^x\) gives
\[
  T_n=\frac{e^{x_n}}{2h}
  \left(e^{2h}-1-2he^h\right).
\]
Using
\[
  e^{2h}-1-2he^h
  =\frac13h^3+\frac13h^4+O(h^5),
\]
we obtain
\[
  T_n=e^{x_n}
  \left(\frac16h^2+\frac16h^3+O(h^4)\right).
\]
Therefore the method has order \(2\), and
\[
  \boxed{C=\frac16}.
\]
\end{solution}

\begin{question}[Constructing a maximal-order implicit formula]\pointmark{10}
Choose \(a,b,c\) so that
\[
  y_{n+2}-y_{n+1}
  =h(af_n+bf_{n+1}+cf_{n+2})
\]
has maximum possible order.  Find its error constant and classify its
zero-stability.
\end{question}

\begin{solution}
The conditions \(c_1=c_2=c_3=0\) are
\[
  a+b+c=1,\qquad
  b+2c=\frac32,\qquad
  b+4c=\frac73.
\]
Solving,
\[
  a=-\frac1{12},\qquad
  b=\frac23,\qquad
  c=\frac5{12}.
\]
Thus
\[
  y_{n+2}-y_{n+1}
  =\frac h{12}(-f_n+8f_{n+1}+5f_{n+2}).
\]
The next coefficient is
\[
  c_4=\frac58-\frac23=-\frac1{24}\neq0.
\]
Since \(\sigma(1)=1\), the method has order \(3\) and error constant
\[
  \boxed{C=-\frac1{24}}.
\]
Finally, \(\rho(z)=z(z-1)\), whose roots are \(0\) and \(1\).
It is strongly zero-stable.
\end{solution}

\begin{question}[Distinct real roots with initial data]\pointmark{5}
Solve
\[
  y_{n+2}-5y_{n+1}+6y_n=0,
  \qquad y_0=1,\quad y_1=4.
\]
\end{question}

\begin{solution}
The auxiliary equation
\[
  \xi^2-5\xi+6=(\xi-2)(\xi-3)=0
\]
has distinct roots \(2,3\).  Thus
\[
  y_n=A2^n+B3^n.
\]
The initial conditions give \(A+B=1\) and \(2A+3B=4\), whence
\(A=-1\), \(B=2\).  Therefore
\[
  \boxed{y_n=-2^n+2\cdot3^n}.
\]
\end{solution}

\begin{question}[Repeated real root with initial data]\pointmark{5}
Solve
\[
  y_{n+2}-6y_{n+1}+9y_n=0,
  \qquad y_0=2,\quad y_1=9.
\]
\end{question}

\begin{solution}
The auxiliary equation is \((\xi-3)^2=0\), so
\[
  y_n=(A+Bn)3^n.
\]
From \(y_0=2\), \(A=2\).  From \(y_1=9\),
\(3(2+B)=9\), so \(B=1\).  Hence
\[
  \boxed{y_n=(n+2)3^n}.
\]
\end{solution}

\begin{question}[Complex roots with initial data]\pointmark{6}
Solve
\[
  y_{n+2}-2y_{n+1}+4y_n=0,
  \qquad y_0=1,\quad y_1=0.
\]
\end{question}

\begin{solution}
The roots are
\[
  1\pm i\sqrt3
  =2e^{\pm i\pi/3}.
\]
Thus
\[
  y_n=2^n\left(A\cos\frac{n\pi}{3}
  +B\sin\frac{n\pi}{3}\right).
\]
The first condition gives \(A=1\).  The second gives
\[
  2\left(\frac12+\frac{\sqrt3}{2}B\right)=0,
\]
so \(B=-1/\sqrt3\).  Therefore
\[
  \boxed{
  y_n=2^n\left(\cos\frac{n\pi}{3}
  -\frac1{\sqrt3}\sin\frac{n\pi}{3}\right)}.
\]
\end{solution}

\begin{question}[Repeated complex-conjugate roots]\pointmark{6}
Find the general real solution of
\[
  (E^2-2E+2)^2y_n=0.
\]
\end{question}

\begin{solution}
The roots \(1\pm i=\sqrt2e^{\pm i\pi/4}\) each have multiplicity two.
The general real solution is
\[
  \boxed{
  y_n=(\sqrt2)^n
  \left[
  (A+Bn)\cos\frac{n\pi}{4}
  +(C+Dn)\sin\frac{n\pi}{4}
  \right]}.
\]
\end{solution}

\begin{question}[Constant and nonresonant exponential forcing]\pointmark{8}
Solve
\[
  y_{n+2}-5y_{n+1}+6y_n=4+4^n,
  \qquad y_0=0,\quad y_1=1.
\]
\end{question}

\begin{solution}
The complementary solution is \(A2^n+B3^n\).  Since
\[
  \phi(1)=2,\qquad \phi(4)=2,
\]
particular solutions for \(4\) and \(4^n\) are respectively
\[
  2,\qquad \frac{4^n}{2}.
\]
Thus
\[
  y_n=A2^n+B3^n+2+\frac{4^n}{2}.
\]
The initial conditions yield
\[
  A+B=-\frac52,\qquad 2A+3B=-3,
\]
so \(A=-9/2\), \(B=2\).  Therefore
\[
  \boxed{
  y_n=-\frac92\,2^n+2\cdot3^n+2+\frac{4^n}{2}}.
\]
\end{solution}

\begin{question}[Simple exponential resonance]\pointmark{8}
Solve
\[
  y_{n+2}-5y_{n+1}+6y_n=2^n,
  \qquad y_0=1,\quad y_1=0.
\]
\end{question}

\begin{solution}
The complementary solution is \(A2^n+B3^n\).  Since \(2\) is a
simple root of \(\phi(\xi)=(\xi-2)(\xi-3)\), use
\[
  y_n^{(p)}
  =\frac{n2^n}{2\phi'(2)}
  =-\frac12n2^n.
\]
Thus
\[
  y_n=A2^n+B3^n-\frac12n2^n.
\]
The initial conditions give \(A+B=1\) and \(2A+3B=1\), so
\(A=2\), \(B=-1\).  Hence
\[
  \boxed{y_n=2\cdot2^n-3^n-\frac12n2^n}.
\]
\end{solution}

\begin{question}[Nonresonant polynomial forcing]\pointmark{7}
Find the general solution of
\[
  y_{n+2}-5y_{n+1}+6y_n=n^2.
\]
\end{question}

\begin{solution}
Because \(\phi(1)=2\neq0\), take
\[
  y_n^{(p)}=an^2+bn+c.
\]
Substitution gives
\[
  2an^2+(-6a+2b)n+(-a-3b+2c)=n^2.
\]
Hence
\[
  a=\frac12,\qquad b=\frac32,\qquad c=\frac52.
\]
Therefore
\[
  \boxed{
  y_n=A2^n+B3^n+\frac12n^2+\frac32n+\frac52}.
\]
\end{solution}

\begin{question}[Nonresonant polynomial times an exponential]\pointmark{9}
Find the general solution of
\[
  y_{n+2}-5y_{n+1}+6y_n=4^n(n^2+1).
\]
\end{question}

\begin{solution}
Put \(y_n^{(p)}=4^nv(n)\), where
\[
  v(n)=an^2+bn+c.
\]
The shift rule gives
\[
  (16E^2-20E+6)v(n)=n^2+1.
\]
Equating coefficients yields
\[
  a=\frac12,\qquad b=-6,\qquad c=\frac{51}{2}.
\]
Adding the complementary solution gives
\[
  \boxed{
  y_n=A2^n+B3^n+
  4^n\left(\frac12n^2-6n+\frac{51}{2}\right)}.
\]
\end{solution}

\begin{question}[Polynomial times a simply resonant exponential]\pointmark{10}
Find the general solution of
\[
  y_{n+2}-5y_{n+1}+6y_n=2^n(n+1).
\]
\end{question}

\begin{solution}
Set \(y_n^{(p)}=2^nv(n)\).  Then
\[
  (4E^2-10E+6)v(n)=n+1.
\]
Because \(2\) is a simple characteristic root, a polynomial \(v\) of
one degree higher than the forcing is required.  Taking
\[
  v(n)=-\frac14n^2-\frac54n
\]
gives the required right-hand side.  Hence
\[
  \boxed{
  y_n=A2^n+B3^n
  -\frac{2^n}{4}(n^2+5n)}.
\]
The arbitrary constant term in \(v\) would only reproduce the
complementary term \(A2^n\).
\end{solution}

\begin{question}[Conversion from \(\Delta\)-form and mixed forcing]\pointmark{10}
Solve
\[
  \Delta^2y_n-3\Delta y_n+2y_n=5^n+n,
  \qquad y_0=0,\quad y_1=1.
\]
\end{question}

\begin{solution}
Since \(\Delta=E-1\),
\[
  \Delta^2-3\Delta+2
  =(E-1)^2-3(E-1)+2=E^2-5E+6.
\]
Thus the complementary solution is \(A2^n+B3^n\).  For the
exponential forcing,
\[
  \phi(5)=25-25+6=6,
\]
so its particular solution is \(5^n/6\).  For the linear forcing,
substitution of \(an+b\) gives
\[
  a=\frac12,\qquad b=\frac34.
\]
Therefore
\[
  y_n=A2^n+B3^n+\frac{5^n}{6}+\frac n2+\frac34.
\]
The initial conditions yield
\[
  A+B=-\frac{11}{12},\qquad
  2A+3B=-\frac{13}{12}.
\]
Hence \(A=-5/3\), \(B=3/4\), and
\[
  \boxed{
  y_n=-\frac53\,2^n+\frac34\,3^n
  +\frac{5^n}{6}+\frac n2+\frac34}.
\]
\end{solution}

\end{document}
